%# -*- coding: utf-8-unix -*-
\chapter{搭建模板编译环境}

\section{Windows 适配版使用说明}

本版本面向 Windows 用户，基于旧版 SJTU-TermPaper 课程论文模板进行适配。
所用文档类为 \texttt{sjtuthesis}，版本为 2017 年 2 月 19 日的 \texttt{v0.9.5}。
2026 年 10 月 28 日，已在本机 Windows 与 TeX Live 2024 环境中完成
XeLaTeX、Biber、XeLaTeX、XeLaTeX 的完整编译验证，成功生成 PDF。

\subsection{安装与编译}

Windows 用户需安装可用的 TeX Live 及模板所需宏包，并确认系统中有
宋体（SimSun）、黑体（SimHei）、楷体（KaiTi）和仿宋（FangSong）。
本版本默认采用 \texttt{fontset=windows}，无需安装 Adobe 中文字体。
安装完成后，确认终端可以调用 \texttt{xelatex} 和 \texttt{biber}。

在主文件 \texttt{thesis.tex} 所在目录打开 PowerShell 或命令提示符，依次执行：

\begin{lstlisting}[basicstyle=\small\ttfamily, numbers=none, breaklines=true]
xelatex -interaction=nonstopmode -halt-on-error -file-line-error thesis.tex
biber thesis
xelatex -interaction=nonstopmode -halt-on-error -file-line-error thesis.tex
xelatex -interaction=nonstopmode -halt-on-error -file-line-error thesis.tex
\end{lstlisting}

编译结果为模板根目录下的 \texttt{thesis.pdf}。
使用图形编辑器时，须将 \texttt{thesis.tex} 设为主文件，选择 XeLaTeX 编译器，
并使用 Biber 处理参考文献。

\subsection{本版本修复的问题}

\begin{enumerate}
	\item \textbf{Adobe 中文字体缺失。}
	      原主文件 \texttt{thesis.tex} 指定 \texttt{fontset=adobe}，
	      在未安装相应字体的 Windows 环境中，会出现
	      \texttt{The font "AdobeSongStd-Light" cannot be found} 错误并中止编译。
	      本版本将默认设置改为 \texttt{fontset=windows}，使用上述 Windows 中文字体。

	\item \textbf{双语题注选项在声明前被使用。}
	      原文档类 \texttt{sjtuthesis.cls} 在声明
	      \texttt{bi-first}、\texttt{bi-second} 选项之前调用它们，
	      在本机宏包环境下触发 \texttt{Package caption Error: bi-first undefined}。
	      本版本先通过 \verb|\DeclareCaptionOption| 声明选项，
	      再执行相应的 \verb|\captionsetup| 设置。

	\item \textbf{嵌套条件缺少结束标记。}
	      原配置文件 \texttt{sjtuthesis.cfg} 在学位类型判断中缺少一个
	      \verb|\fi|，导致编译结束时提示条件未闭合。
	      本版本已补齐该标记，使嵌套条件完整闭合。
\end{enumerate}

字体配置是针对 Windows 的默认适配；题注选项顺序与条件闭合问题属于旧版模板代码问题。
此外，页脚总页数改用已加载的 \texttt{lastpage} 宏包提供的
\texttt{LastPage} 标签，使增加附录内容后仍能正确统计总页数。
当前示例正文未引用参考文献，因此参考文献为空或 Biber 提示无引用均属预期情况。
字体尺寸替代、行宽超出等排版警告不等同于编译失败。

\subsection{版本来源与适用范围}

本次修复针对本目录中的旧版课程论文模板。
当前上游项目为 \href{https://github.com/sjtug/SJTUThesis}{\texttt{sjtug/SJTUThesis}}，
其文件结构与本旧版模板不同；本次验证不表示当前上游存在相同问题，
也不保证可将这些修改直接套用到其最新版本。
后续 Fork 当前上游仓库时，应按其结构另行整合并验证。

本版本默认面向 Windows。在 macOS、Linux 或在线平台上使用时，
应选择当地可用的 CTeX 字体配置，并重新进行编译验证。
以下保留原模板中的其他平台安装说明作为历史参考；
其中的发行版年份及命令并非本次验证环境。

\section{安装TeX发行版}

\subsection{Mac OS X}

Mac用户可以从MacTeX主页\footnote{\url{https://tug.org/mactex/}}下载MacTeX 2015。
也可以通过brew包管理器\footnote{\url{http://caskroom.io}}安装MacTeX 2015。

\begin{lstlisting}[basicstyle=\small\ttfamily, numbers=none]
brew cask install mactex
\end{lstlisting}

\subsection{Linux}

建议Linux用户使用TeXLive主页\footnote{\url{https://www.tug.org/texlive/}}的脚本来安装TeXLive 2015。
以下命令将把TeXLive发行版安装到当前用户的家目录下。
若计划安装一个供系统上所有用户使用的TeXLive，请使用root账户操作。

\begin{lstlisting}[basicstyle=\small\ttfamily, numbers=none]
wget http://mirror.ctan.org/systems/texlive/tlnet/install-tl-unx.tar.gz
tar xzvpf install-tl-unx.tar.gz
cd install-tl-20150411/
./install-tl
\end{lstlisting}

\section{安装中文字体}

\subsection{Mac OS X、Deepin}

Mac和Deepin用户双击字体文件即可安装字体。

\clearpage
\subsection{RedHat/CentOS用户}

RedHat/CentOS用户请先将字体文件复制到字体目录下，调用fc-cache刷新缓存后即可在TeXLive中使用新字体。

\begin{lstlisting}[basicstyle=\small\ttfamily, numbers=none]
mkdir ~/.fonts
cp *.ttf ~/.fonts				# 当前用户可用新字体
cp *.ttf /usr/share/fonts/local/	# 所有用户可以使用新字体
fc-cache -f
\label{last}
\end{lstlisting}

